\input amstex
\documentstyle{amsppt}
\magnification 1200
\pageno 53
\hsize 6.5 true in
\vsize 8.7 true in
\define\zzz{\Bbb Z}
\define\qqq{\Bbb Q}
\define\ccc{\Bbb C} 
\define\aaa{\Bbb A} 
\define\rrr{\Bbb R}
\define\eps{\epsilon}
\define\okk{{\Cal O}_K}
\baselineskip 16pt
\document
\heading
{\bf Appendix: Background on $G_2$}
\endheading
\par
\smallskip
The first four sections of this appendix have been taken from an 
unpublished joint paper with D. Bump. 
\par
The purpose of this appendix is to collect some foundational material on
the exceptional group $G_2$, then to use this material to 
compute the cuspidal endoscopic groups of 
$G_2$, the groups $H^1(\Gamma, T)$, and the $F$-classes in the
stable 2- and 3-regular unipotent orbit.
\par
In the first few sections, we shall be concerned with the phenomenon 
known as ``triality'',
which yields an embedding of $G_2$ into the double cover Spin(8) of
the split form of SO(8). We shall
determine the triality map explicitly at the root level, and also use it to 
give an explicit Chevalley basis for the Lie algebra ${\frak g}_2$, 
which should be useful in other calculations.
\newline
Acknowledgement: We thank Becci Davies for coding the root diagrams 
(in AMS--TEX).
\vskip .3in
\subheading{\bf Background on G$_2$ and SO(8)}
\vskip .1in
\subheading{A.1. SO(8)}
\vskip .1in
Let J denote the $8\times 8$ matrix with 1's on the skew--diagonal and 
0's elsewhere:
$$
J := \pmatrix 0_{\quad} &&{}_{_{\dsize\cdot}\quad}1 \\
   \hfil&{}^{\dsize\cdot}&\\
   1^{\quad^{\dsize\cdot}}& &{}_{_{\quad}}0
   \endpmatrix.
$$
Define the $8\times 8$ {\bf special orthogonal group} 
by
$$
SO(8) := \{A \in GL(8) \ |\ A\cdot J\cdot ^{t}A = J\}^o,
$$
so that $SO(8)$ is semi--simple, connected, neither simply connected
nor adjoint, non--compact, and split over $\qqq$. Its maximal $\qqq$-split
torus consists of all matrices of the form
$$
\pmatrix A&0\\0&B \endpmatrix,\ \text{with }
A:=\tsize{\pmatrix a_1&\hfil&\hfil&\\
\hfil&a_2&\hfil&\\
\hfil&\hfil&a_3&\\
\hfil&\hfil&\hfil&a_4 \endpmatrix},\ 
B:= \tsize{\pmatrix a_4^{-1}&\hfil&\hfil&\\
\hfil&a_3^{-1}&\hfil&\\
\hfil&\hfil&a_2^{-1}&\\
\hfil&\hfil&\hfil&a_1^{-1} \endpmatrix}.\tag A.1
$$
The compact real form of $SO(8)$ is obtained, of course, by replacing
J by the $8\times 8$ identity matrix. 
\par
The simply--connected covering group of $SO(8)$ is
$SO(8)_{sc} =$ Spin(8), which has center $\zzz /2\zzz \times \zzz /2\zzz$,
and since $SO(8)$ has center $\zzz /2\zzz$, Spin(8) is a 2-fold cover
of $SO(8)$ and a 4-fold cover of $PSO(8)$. The outer automorphisms
of Spin(8) are in one--to--one correspondence with the automorphisms
of the Dynkin diagram for $SO(8)$:
$$
\matrix
\hfil&\hfil&\hfil&\quad\bullet \\
\hfil&\hfil&\hfil\\ 
\hfil&\hfil&\hfil\\ 
\bullet\quad\quad&\hfil&\bullet
\ &\hfil \\
\hfil&\hfil&\hfil\\ 
\hfil&\hfil&\hfil\\ 
\hfil&\hfil&\hfil&\quad\bullet\endmatrix
$$
The Lie algebra
$$
so(8) := \{A \in M_8\ |\ AJ + J^{t}A = 0\}
$$
may be expressed as the Lie algebra of $8\times 8$ matrices of the
form
$$
A =  {\pmatrix t_{1}&a_{12}&a_{13}&a_{14}&a_{15}&a_{16}&a_{17}&0\\
a_{21}&t_2&a_{23}&a_{24}&a_{25}&a_{26}&0&-a_{17}\\
a_{31}&a_{32}&t_3&a_{34}&a_{35}&0&-a_{26}&-a_{16}\\
\vdots&\vdots &\vdots &\vdots &\vdots &\vdots &\vdots &\vdots\\
0&-a_{71}&-a_{61}&-a_{51}&-a_{41}&-a_{31}&-a_{21}&-t_1\\
\endpmatrix}.\tag A.2
$$
If $i \not= j$ then define $E_{ij}$ to be the matrix A as above with
$a_{ij} = 1$ and all other entries = 0. Define $E_{ii}$ to be the
matrix A as above with $t_i = 1$ and all other entries = 0. Then
$$
\split
&\big\{ E_{11},\,  E_{22},\, E_{33},\, E_{44},\\
&\quad E_{12},\, E_{13},\dots ,E_{17},
       E_{23},\dots ,E_{26},
       E_{34},\, E_{35},\\
&\quad {^{t}E_{12}},\, {^{t}E_{13}},\dots ,{^{t}E_{17}},
       {^{t}E_{23}},\dots ,{^{t}E_{26}},
       {^{t}E_{34}},{^{t}E_{35}} \big\}
\endsplit\tag A.3
$$
forms a basis
for the 28-dimensional vector space $so(8)$. 
\vskip .3in
\subheading{A.2. The Cayley algebra and $G_2$}
\vskip .1in
The {\bf split Cayley algebra} is the $2\times 2$ matrix algebra
$$
C:=\Big\{\tsize{\pmatrix \alpha&a\\b&\beta \endpmatrix}\ | \ 
\alpha ,\beta \in \qqq,\ a,b \in V \Big\},\tag A.4
$$
where $V \cong \qqq^3$ is the cross product algebra
$$
\split
V:= \{ xi+yj+zk\ | \ x,y,z \in \qqq,
&\ i\times i = j\times j = k\times k = 0,\\
&\quad\quad
\ i\times j = k,\ j\times k = i,\ k\times j = i,\ j\times i = -k,
\ \text{etc.}\}.
\endsplit
$$
Define 
$$
G_2(F) :=\text{Aut}_F(C\otimes F),\tag A.5
$$
where F is any field containing $\qqq$. Since Der(C) $\cong$ Lie(Aut C),
it follows that the Lie algebra ${\frak g}_2$ may be 
identified with Der(C). For $G_2$, there are 2 real forms:  one
being the non-compact $\qqq$--split from above, the other being a
compact form associated to a division algebra (see\cite{Jac, pp. 142-147}
or \cite{T}). We have
$$
\text{dim}_{_{\qqq}}C = 8,\ \text{dim}_{_{\qqq}}{\frak g}_2 = 14,
$$
and there is a trace bilinear form which induces an embedding
$$
G_2 = \text{Aut}(C) \hookrightarrow SO(8).
$$
Here SO(8) is regarded as the special orthogonal group associated to
the trace form $trace(\beta)$ on C, where
$$
\beta(x,y)={1\over 2}(x\cdot y^*+y^*\cdot x),
\ \ \ \ \beta(x,x)=N(x),
$$
and where $N$ denotes the norm on $C$. Regarding $G_2$ as a subgroup of
$SO(8)$, we have the following 
\proclaim{Triality theorem A.6(\cite{Jac})} Given any orthogonal map
$$
A_0 : C \rightarrow C,
$$
there are orthogonal maps $A_1,\ A_2$ such that, for all $x,\,y \in C$,
$$
A_0(xy) = A_1(x)A_2(y).
$$
Moreover, the only other orthogonal maps satisfying this property are
$-A_1,\ -A_2$, so 
$$
A_0 \longmapsto (A_1,A_2)
$$
is well--defined, up to signs.
\endproclaim
\demo{proof} We will verify that our form of triality is the global 
equivalent of Jacobson's form \cite{Jac, p. 8}. Let
$$
\split
&A_0 = 1+\delta B_0,\\
&A_1 = 1+\delta B_1,\\
&A_2 = 1+\delta B_2,
\endsplit
$$
in $SO(8,\qqq )$ (regarded as a group of orthogonal maps with respect to the
trace form on C), where $\delta^2 = 0$ and $B_i \in so(8,\qqq )$. Thus 
$A_0,\ A_1,\ A_2$ may be regarded as lying in the tangent space to SO(8)
at the identity (cf. \cite{Bor, chapter 0}). Then
$A_0(xy) = A_1(x)A_2(y)$ is equivalent to 
$$
B_0(xy) = B_1(x)y+ xB_2(y).
$$
This is Jacobson's expression.
\enddemo
\vskip .1in\noindent
\proclaim{Corollary A.7} The map
$$
\split
SO(8,F) &\rightarrow PSO(8,F),\\
\quad A_0 &\longmapsto \quad \pm A_1,
\endsplit
$$
defines a projective representation of SO(8), i.e., belongs to
PSO(8). In particular, it induces an outer automorphism of degree 3
(the ``triality'' automorphism)
corresponding to a rotation of the Dynkin diagram of $PSO(8,F)$.
\endproclaim
\demo{Remarks}(A.8) If $A_0 \in G_2 = \text{Aut}(C)$ then triality sends 
$A_0 \longmapsto A_0$. Conversely, if triality sends $A_0 \longmapsto A_0$
then $A_0 \in G_2$, so $G_2$ is the subgroup fixed by triality. In
general, triality sends (up to signs)
$$
A_0 \longmapsto A_1,\ A_1 \longmapsto A_2,\ A_2 \longmapsto A_0.
$$
This way we can explicitly construct an outer automorphism of $SO(8)$. 
\par 
(A.9) Triality can be described explicitly on the standard Borel
subgroup B of PSO(8): triality cyclically permutes the ``root groups''
$U_{\alpha}$, defined below, in the same way as it permutes the simple roots 
via its action on the Dynkin diagram; it also permutes the embeddings
$PSL(2) \hookrightarrow PSO(8)$ associated to the simple roots. The
images of the maximal split torus $T \subset PSL(2)$ embedded in $PSO(8)$
generate the maximal split torus in $PSO(8)$ and, of course, the simple
root groups generate the unipotent radical of B. Putting these all together
explicitly gives the action of triality on B.
\enddemo
\vskip .3in
\subheading{A.3. Root systems for $G_2$ and $SO(8)$}
\vskip .1in
Since the symmetric group $S_3$ acts on the Dynkin diagram of $SO(8)$ 
by triality,
it acts on its 24 roots. Of these, 6 roots remain fixed under 
this action, and 18 are permuted cyclically in groups of three. In
particular, the roots corresponding to
$$
\{E_{17}\},\ \ \{E_{16}\},\ \ \{E_{23}\},
$$
are fixed and the triples of roots corresponding to
$$
\{E_{14},\ E_{15},\ E_{26}\},\ \{E_{13},\ E_{24},\ E_{25}\},
\ \{E_{12},\ E_{34},\ E_{35}\},
$$
are each permuted by the action of $S_3$. More explicitly, let us 
represent a basis of $so(8)$ by
$$
\eps_1 - \eps_2,\  \eps_2 - \eps_3,\ \eps_3 - \eps_4,\ \eps_3 + \eps_4,
$$
where 
$$
\eps_i : \tsize{\pmatrix A&0\\0&B\endpmatrix}
\longmapsto a_i, \tag A.10
$$
in the notation of (A.1), $1 \leq i \leq 4$.
Here $E_{12}$ corresponds to $\eps_1 - \eps_2,\ E_{13}$ to
$\eps_1 - \eps_2 + \eps_2 - \eps_3 = \eps_1 - \eps_3$, etc. Since
$\eps_2 - \eps_3$ is connected to the others (since its Killing form inner
product with them is non--zero), it corresponds to the center of the
Dynkin diagram
$$
\matrix
\hfil&\hfil&\hfil&\quad\bullet\text{$\eps_1-\eps_2$} \\
\hfil&\hfil&\hfil\\ 
\hfil&\hfil&\hfil\\ 
\eps_3-\eps_4\ \bullet\quad\quad&\hfil&\bullet
\ \text{$\eps_2-\eps_3$}&\hfil \\
\hfil&\hfil&\hfil\\ 
\hfil&\hfil&\hfil\\ 
\hfil&\hfil&\hfil&\quad\bullet\text{$\eps_3+\eps_4$}
\endmatrix
$$
Therefore, under the action of $S_3$, the triality automorphism of
so(8) induces a linear map $\phi$ on the roots such that
$$
\phi (\eps_1 - \eps_2) = \eps_3 - \eps_4,
\ \phi (\eps_3 - \eps_4) = \eps_3 + \eps_4,
\ \phi (\eps_3 + \eps_4) = \eps_1 - \eps_2,
\ \phi (\eps_2 - \eps_3) = \eps_2 - \eps_3.\tag A.11
$$
Thus
$$
\split
&\phi (\eps_1) = {1\over{2}}(\eps_1 + \eps_2 + \eps_3 - \eps_4),\\
&\phi (\eps_2) = {1\over{2}}(\eps_1 + \eps_2 - \eps_3 + \eps_4),\\
&\phi (\eps_3) = {1\over{2}}(\eps_1 - \eps_2 + \eps_3 + \eps_4),\\
&\phi (\eps_4) = {1\over{2}}(\eps_1 - \eps_2 - \eps_3 - \eps_4).
\endsplit
$$
Call the induced automorphism on the Lie algebra so(8), $\phi^*$. As
we've already observed, the Lie subalgebra
$$
\{X \in so(8)\ |\ \phi^*(X) = X\}
$$
is the Lie algebra ${\frak g}_2$. Identifying its Lie
group with a subgroup of $SO(8)$, we find that its maximal
$\qqq-$split torus is of the form
$$
\split
T':= \Big\{\pmatrix A&0\\0&B \endpmatrix,\ |\ 
&A:={\left(\smallmatrix a_1&\hfil&\hfil&\\
\hfil&a_2&\hfil&\\
\hfil&\hfil&a_3&\\
\hfil&\hfil&\hfil&1 \endsmallmatrix\right)},\ 
B:= {\left(\smallmatrix 1&\hfil&\hfil&\\
\hfil&a_3^{-1}&\hfil&\\
\hfil&\hfil&a_2^{-1}&\\
\hfil&\hfil&\hfil&a_1^{-1} \endsmallmatrix\right)}\\
&\quad\quad\quad\quad\text{and } a_1 = a_2a_3\Big\}.
\endsplit\tag A.12
$$ 
We will identify $T'$ with the maximal split torus of 
$SL(3)$ by means of the map
$$
\pmatrix A&0\\0&B \endpmatrix \in T' \longmapsto
\text{diag}(a_2,a_3,a_1^{-1}),
$$                    
with the maximal split torus of 
$PGL(3)$ by means of the map
$$
\pmatrix A&0\\0&B \endpmatrix \in T' \longmapsto
\text{diag}(a_1a_2,a_1a_3,1),\tag A.13
$$                    
with the maximal split torus of 
$PGL(2)\times SL(2)$ by means of the map
$$
\pmatrix A&0\\0&B \endpmatrix \in T' \longmapsto
\text{diag}(a_3,1)\times diag(a_1a_2,1/a_1a_2),
\tag A.13'
$$                    
and with the maximal split torus of 
$(SL(2)\times SL(2))/\{\pm 1\}$ by means of the map
$$
\pmatrix A&0\\0&B \endpmatrix \in T' \longmapsto
\text{diag}(a_1,1/a_1)\times diag(a_2,1/a_2).
\tag A.13''
$$                    
\vskip .3in
\subheading{\bf A.4. The Chevalley basis}
\vskip .1in
The embedding of ${\frak g}_2$ into so(8) gives a convenient framework 
for doing calculations with explicit matrices. We now write down
the Chevalley basis of $G_2$. The notation for the root system is given
by the following diagram:
$$
\tsize{\matrix \hfil &\hfil&\hfil&\hfil &\overset\text{3$\alpha + 2\beta$}\to\bullet&\hfil&\hfil\\
\hfil&\hfil&\hfil&\hfil&\hfil\\
\hfil&\hfil&\hfil&\hfil&\hfil\\
\hfil&\hfil&\hfil&\hfil&\hfil\\
\overset\text{\ \ \ \ $\beta$\ \ \ \ 
}\to\bullet&\hfil&\hfil&\overset\text{$\alpha +\beta\ \ 
\quad$}\to\bullet&\hfil
    &\overset\text{$\ \ \quad2\alpha + 
\beta$}\to\bullet&\hfil&\hfil&\overset\text{$3\alpha 
+\beta$}\to\bullet\\
\hfil&\hfil&\hfil&\hfil&\hfil\\
\hfil&\hfil&\hfil&\hfil&\hfil\\
\hfil&\hfil&\hfil&\hfil&\hfil\\
\hfil&\hfil&\hfil&\hfil&\hfil\\
\hfil&\hfil&\hfil&\hfil&\hfil\\
\hfil&\hfil&\hfil&\hfil&\hfil\\
  \hfil&\bullet&\hfil&\hfil&\cdot&\hfil&\hfil&\overset\text{$\alpha$}\to\bullet\\
\hfil&\hfil&\hfil&\hfil&\hfil\\
\hfil&\hfil&\hfil&\hfil&\hfil\\
\hfil&\hfil&\hfil&\hfil&\hfil\\
\hfil&\hfil&\hfil&\hfil&\hfil\\
\hfil&\hfil&\hfil&\hfil&\hfil\\
\hfil&\hfil&\hfil&\hfil&\hfil\\
    \bullet&\hfil&\hfil&\bullet\ \ \quad&\hfil&\ \ \quad\bullet&\hfil&\hfil&\bullet\\
\hfil&\hfil&\hfil&\hfil&\hfil\\
\hfil&\hfil&\hfil&\hfil&\hfil\\
\hfil&\hfil&\hfil&\hfil&\hfil\\
    \hfil&\hfil&\hfil&\hfil&\bullet&\hfil&\hfil&\hfil&\hfil
    \endmatrix}
$$
\centerline{the root system for $G_2$}
\vskip .2in
$$
X_{\alpha} :=
\pmatrix 0&1&\ &\ &\ &\ &\ &\ \\
         \ &0&\ &\ &\ &\ &\ &\ \\
         \ &\ &0&1&1&\ &\ &\ \\
         \ &\ &\ &0&0&-1&\ &\ \\
         \ &\ &\ &\ &0&-1&\ &\ \\
         \ &\ &\ &\ &\ &0&\ &\ \\
         \ &\ &\ &\ &\ &\ &0&-1\\
         \ &\ &\ &\ &\ &\ &\ &0
         \endpmatrix,\ \ 
h_{\alpha} := \pmatrix 1&\ &\ &\ &\ &\ &\ &\ \\
                         \ &-1&\ &\ &\ &\ &\ &\ \\
                         \ &\ &2&\ &\ &\ &\ &\ \\
                         \ &\ &\ &0&\ &\ &\ &\ \\
                         \ &\ &\ &\ &0&\ &\ &\ \\
                         \ &\ &\ &\ &\ &-2&\ &\ \\
                         \ &\ &\ &\ &\ &\ &1&\ \\
                         \ &\ &\ &\ &\ &\ &\ &-1
                         \endpmatrix.
$$
(Note that $X_{\alpha} = E_{12} + E_{34} + E_{35}$ and 
$\alpha\pmatrix A&\ \\ \ &B\endpmatrix =a_3$, where $A,B$ are as
in (A.12).)
$$
X_{\beta} :=
\pmatrix 0&\ &\ &\ &\ &\ &\ &\ \\
         \ &0&1&\ &\ &\ &\ &\ \\
         \ &\ &0&\ &\ &\ &\ &\ \\
         \ &\ &\ &0&\ &\ &\ &\ \\
         \ &\ &\ &\ &0&\ &\ &\ \\
         \ &\ &\ &\ &\ &0&-1&\ \\
         \ &\ &\ &\ &\ &\ &0&\\
         \ &\ &\ &\ &\ &\ &\ &0
         \endpmatrix,\ \ 
h_{\beta} := \pmatrix 0&\ &\ &\ &\ &\ &\ &\ \\
                         \ &1&\ &\ &\ &\ &\ &\ \\
                         \ &\ &-1&\ &\ &\ &\ &\ \\
                         \ &\ &\ &0&\ &\ &\ &\ \\
                         \ &\ &\ &\ &0&\ &\ &\ \\
                         \ &\ &\ &\ &\ &1&\ &\ \\
                         \ &\ &\ &\ &\ &\ &-1&\ \\
                         \ &\ &\ &\ &\ &\ &\ &0
                         \endpmatrix.
$$
(Note that $X_{\beta} = E_{23}$ and 
$\beta\pmatrix A&\ \\ \ &B\endpmatrix =a_2a_3^{-1}$, where $A,B$ are as
in (A.12).) For $\sigma = \alpha+\beta$:
$$
X_{\sigma} :=
\pmatrix 0&\ &-1&\ &\ &\ &\ &\ \\
         \ &0&\ &1&1&\ &\ &\ \\
         \ &\ &0&\ &\ &\ &\ &\ \\
         \ &\ &\ &0&\ &\ &-1&\ \\
         \ &\ &\ &\ &0&\ &-1&\ \\
         \ &\ &\ &\ &\ &0&\ &1\\
         \ &\ &\ &\ &\ &\ &0&\ \\
         \ &\ &\ &\ &\ &\ &\ &0
         \endpmatrix,\ \ 
h_{\sigma} := \pmatrix 1&\ &\ &\ &\ &\ &\ &\ \\
                         \ &2&\ &\ &\ &\ &\ &\ \\
                         \ &\ &-1&\ &\ &\ &\ &\ \\
                         \ &\ &\ &0&\ &\ &\ &\ \\
                         \ &\ &\ &\ &0&\ &\ &\ \\
                         \ &\ &\ &\ &\ &1&\ &\ \\
                         \ &\ &\ &\ &\ &\ &-2&\ \\
                         \ &\ &\ &\ &\ &\ &\ &-1
                         \endpmatrix.
$$
(Note that $X_{\alpha+\beta} = E_{13}+E_{24}+E_{25}.$) For
$\sigma = 2\alpha+\beta$:
$$
X_{\sigma} :=
\pmatrix 0&\ &\ &1&1&\ &\ &\ \\
         \ &0&\ &\ &\ &1&\ &\ \\
         \ &\ &0&\ &\ &\ &-1&\ \\
         \ &\ &\ &0&\ &\ &\ &-1\\
         \ &\ &\ &\ &0&\ &\ &-1\\
         \ &\ &\ &\ &\ &0&\ &\ \\
         \ &\ &\ &\ &\ &\ &0&\ \\
         \ &\ &\ &\ &\ &\ &\ &0
         \endpmatrix,\ \ 
h_{\sigma} := \pmatrix 2&\ &\ &\ &\ &\ &\ &\ \\
                         \ &1&\ &\ &\ &\ &\ &\ \\
                         \ &\ &1&\ &\ &\ &\ &\ \\
                         \ &\ &\ &0&\ &\ &\ &\ \\
                         \ &\ &\ &\ &0&\ &\ &\ \\
                         \ &\ &\ &\ &\ &-1&\ &\ \\
                         \ &\ &\ &\ &\ &\ &-1&\ \\
                         \ &\ &\ &\ &\ &\ &\ &-2
                         \endpmatrix.
$$
(Note that $X_{2\alpha + \beta} = E_{14}+E_{15}+E_{26}$.) For
$\sigma = 3\alpha+\beta$:
$$
X_{\sigma} :=
\pmatrix 0&\ &\ &\ &\ &1&\ &\ \\
         \ &0&\ &\ &\ &0&\ &\ \\
         \ &\ &0&\ &\ &0&0&-1\\
         \ &\ &\ &0&\ &\ &\ &\ \\
         \ &\ &\ &\ &0&\ &\ &\ \\
         \ &\ &\ &\ &\ &0&\ &\ \\
         \ &\ &\ &\ &\ &\ &0&\ \\
         \ &\ &\ &\ &\ &\ &\ &0
         \endpmatrix,\ \ 
h_{\sigma} := \pmatrix 1&\ &\ &\ &\ &\ &\ &\ \\
                         \ &0&\ &\ &\ &\ &\ &\ \\
                         \ &\ &1&\ &\ &\ &\ &\ \\
                         \ &\ &\ &0&\ &\ &\ &\ \\
                         \ &\ &\ &\ &0&\ &\ &\ \\
                         \ &\ &\ &\ &\ &-1&\ &\ \\
                         \ &\ &\ &\ &\ &\ &0&\ \\
                         \ &\ &\ &\ &\ &\ &\ &-1
                         \endpmatrix.
$$
(Note that $X_{3\alpha + \beta} = E_{16}$.) For $\sigma =
3\alpha+2\beta$:
$$
X_{\sigma} :=
\pmatrix 0&\ &\ &\ &\ &\ &1&\ \\
         \ &0&\ &\ &\ &\ &0&-1\\
         \ &\ &0&\ &\ &\ &\ &\ \\
         \ &\ &\ &0&\ &\ &\ &\ \\
         \ &\ &\ &\ &0&\ &\ &\ \\
         \ &\ &\ &\ &\ &0&\ &\ \\
         \ &\ &\ &\ &\ &\ &0&\ \\
         \ &\ &\ &\ &\ &\ &\ &0
         \endpmatrix,\ \ 
h_{\sigma} := \pmatrix 1&\ &\ &\ &\ &\ &\ &\ \\
                         \ &1&\ &\ &\ &\ &\ &\ \\
                         \ &\ &0&\ &\ &\ &\ &\ \\
                         \ &\ &\ &0&\ &\ &\ &\ \\
                         \ &\ &\ &\ &0&\ &\ &\ \\
                         \ &\ &\ &\ &\ &0&\ &\ \\
                         \ &\ &\ &\ &\ &\ &-1&\ \\
                         \ &\ &\ &\ &\ &\ &\ &-1
                         \endpmatrix.
$$                                   
(Note that $X_{3\alpha+2\beta} = E_{17}$.)
\par
There is a canonical embedding $G_2 \hookrightarrow SO(8)$, given by
the natural faithful representation associated to the Cayley algebra.
In fact, each one--parameter unipotent subgroup $U_{\sigma}$ associated to a
short root $\sigma$ of $G_2$ is contained in a three dimensional unipotent 
subgroup of SO(8) generated by three (uniquely determined) roots of 
$D_4$ which are permuted by $S_3$. 
The possibilities for this embedding of unipotents can be tabulated
as follows:
$$
\split
&U_{\alpha} \subset G_2 \longleftrightarrow 
U := \exp[E_{12}F+E_{34}F+E_{35}F],\\
&U_{2\alpha +\beta} \subset G_2 \longleftrightarrow 
U := \exp[E_{14}F+ E_{15}F+E_{26}F],\\
&U_{\alpha+\beta} \subset G_2 \longleftrightarrow 
U := \exp[E_{13}F+E_{24}F+E_{25}F].
\endsplit\tag A.14
$$
For the long roots of $G_2$, we have
$$
\split
&U_{\beta} \subset G_2 \longleftrightarrow 
U := \exp[E_{23}F],\\
&U_{3\alpha +2\beta} \subset G_2 \longleftrightarrow 
U := \exp[E_{17}F],\\
&U_{3\alpha+\beta} \subset G_2 \longleftrightarrow 
U := \exp[E_{16}F].
\endsplit\tag A.15
$$
In other words, under the action of $S_3$ on the roots of 
$D_4$, the six stable roots correspond to the long roots of $G_2$ 
and the other 18, permuted in orbits of 3, correspond to the short 
roots of $ G_2$.
\vskip .1in
\subheading{A.5. Endoscopic groups of $G_2$}
\vskip .1in
The calculation of the cuspidal endoscopic groups of $G_2$ is as
follows. Motivated by facts
about root systems of centralizers of semi-simple elements
\cite{Spal, Ch III, \S 12}, let $\Delta:=\{\alpha,\beta\}$,
$\Delta_1:=\{\alpha,3\alpha+2\beta\}$,
$\Delta_2:=\{\beta,3\alpha+\beta\}$,
$\Delta^{\vee}:=\{\alpha^{\vee},\beta^{\vee}\}$,
$\Delta_1^{\vee}:=\{\alpha^{\vee},(3\alpha+2\beta)^{\vee}\}$,
and $\Delta_2^{\vee}:=\{\beta^{\vee},(3\alpha+\beta)^{\vee}\}$,
where $\sigma^{\vee}=2\sigma/(\sigma,\sigma)$.
Representing the root $\alpha$ by the vector
$(1,0)$ and $\beta$ by $(-{3\over 2},{\sqrt{3}\over 2})$,
we find that the root lattice $Q$ is given by
$$
Q={\Bbb Z}[\Delta]={\Bbb Z}[(1,0),(-{3\over 2},{\sqrt{3}\over 2})].
$$
We have $\alpha^{\vee}=(2,0)$ and $\beta^{\vee}=(-1,{1\over \sqrt{3}})$.
Since $G_2$ is adjoint, its character lattice $X$ is equal to its
root lattice: $X=Q$. $G_2$ is also simply connected, so the weight lattice
$$
\split
P
&=\{x\in {\Bbb R}^2\ |\ (x,\Sigma^{\vee})\subset {\Bbb Z}\}\\
&=\{(a,b)\ |\ (a,b)\cdot (2,0)\in {\Bbb Z},\  
(a,b)\cdot (-1,{1\over \sqrt{3}})\in {\Bbb Z}\}\\
&={\Bbb Z}[(1/2,0),(0,{\sqrt{3}\over 2})]
\endsplit
$$ 
is also equal to the character lattice. Similarly,
$Q^{\vee}={\Bbb Z}[(2,0),(-1,{1\over \sqrt{3}})]=P^{\vee}$.
We have $Q_1={\Bbb Z}[(1,0),(0,\sqrt{3})]$,
$P_1={\Bbb Z}[(1/2,0),(0,\sqrt{3}/2)]$,
$Q_1^{\vee}={\Bbb Z}[(2,0),(0,2/\sqrt{3})]$,
and $P_1^{\vee}={\Bbb Z}[(1,0),(0,1/\sqrt{3})]$.
Fix a maximal split torus $^LT^o$ of $^LG_2^o$ shared by $^LH^o$
and identify $X$ with the lattice of characters of $^LT^o$. Take
$^LH^o$ to be of type $A_1\times A_1$, with root system generated by
$\beta^{\vee}$ (which is now short) and $(3\alpha+2\beta)^{\vee}$
(which is not equal to $3\beta^{\vee}+2\alpha^{\vee}$). 
The character lattice is equal to $X_1^{\vee}=X^{\vee}=Q^{\vee}$ 
but its root lattice $Q_1^{\vee}$ is strictly contained in $Q^{\vee}$ and
$Q^{\vee}=X_1^{\vee}$ strictly contains $P^{\vee}$. Therefore, $^LH^0$
is neither
simply connected nor adjoint if $^LH^o$ is of type $A_1\times A_1$.
This forces either $H=PGL(2)\times SL(2)$ or
$H=(SL(2)\times SL(2))/\{\pm 1\}$. Following a suggestion of T. Hales, we now 
show that the former possibility may be discarded. The center
$Z_{^LH^o}$ is of order two: let $z\in Z_{^LH^o}$ and note that
$\alpha^{\vee}(z)=\zeta^{\vee}(z)=1$ (the dual roots acting on the maximal 
split torus of the dual group); these imply that 
$\beta^{\vee}(z)=\pm 1$, which means that
$$
Z_{^LH^o}({\Bbb C})=\{diag(\epsilon,\epsilon,1,1,1,1,\epsilon,\epsilon)\ |\ 
\text{ where }\epsilon =\pm 1\}.
$$
The image of this in $PGL(2)\times SL(2)$, under the embedding in
(A.13) is trivial, so $H=(SL(2)\times SL(2))/\{\pm 1\}$.
Similarly, using $\Delta_2$ in place of $\Delta_1$, 
one finds that $^LH^o\ \subset\ ^LG_2^o$ 
is simply connected if $^LH^o$ is of type $A_2$. (This fact
can also be found in \cite{SpSt, ch II, 5.3-5.5}.)
Therefore $H=PGL(3)$ and $H=(SL(2)\times SL(2))/\{\pm 1\}$, as desired.
\vskip .1in
\subheading{A.6. Tori of $G_2$}
\vskip .1in
The stable conjugacy classes of Cartans in a quasi-split group $G$
are in 1-1 correspondence with the cohomology classes in 
$H^1(\Gamma, W)$, where $W$ denotes the absolute Weyl group of
a fixed Cartan $T$ of $G$ defined over $F$ \cite{Ro, \S 4, Lemma 3}.
One may also determine the Cartans from the $\Gamma$-module
structure on their lattice of characters \cite{Sat}.
\par
Let $T$ be a Cartan of $G_2$ defined over $F$ which splits over a Galois
extension $K$. The Galois
group $\Gamma_K:=Gal(K/F)$ acts on the character lattice 
$X=X(T)$, as does the Weyl group $W$. Recall 
that $T$ is determined (up to $F$-isomorphism)
from the $\Gamma_K$-module $X$. The $\Gamma_K$-invariant submodules
$$
X^{\Gamma_K}:=\{\chi\in X\ |\ \chi^\sigma =\chi,\ \forall\sigma\in\Gamma_K\},
$$
and
$$
X_{\Gamma_K}:=\{\chi\in X\ |\ \sum_{\sigma\in\Gamma_K}\chi^\sigma =0\},
$$
satisfy the property that $X/X^{\Gamma_K}$ and $X/X_{\Gamma_K}$ are
torsion free. Let $T_0$ denote the annihilator of
$X^{\Gamma_K}$ in $T$ and let $A$ denote the annihilator of 
$X_{\Gamma_K}$ in $T$.
Then $T_0\cap A$ is finite, $T=T_0A$, $A$ is the largest $F$-split
subtorus of $T$, and $T_0$ is the largest $F$-anisotropic subtorus of $T$. 
Let $\Sigma_{\Gamma_K}\subset \Sigma$ denote the closed
subsystem defined by $\Sigma_{\Gamma_K}:=X_{\Gamma_K}\cap \Sigma$, let
$\pi:X\rightarrow X/X_{\Gamma_K} \cong X(A)$, and let $\overline{\Sigma}:=
\pi( \Sigma-\Sigma_{\Gamma_K})$ denote the set of restricted roots 
relative to $A$. By Tate-Nakayama \cite{L', p. 702},
the character $\kappa$ of ${\frak D}(T/F)\subset H^1(\Gamma_K,T)$
may be regarded as a character of the ``cocycle sublattice'' 
$X_{\Gamma_K}^{\vee}$, trivial on the ``coboundary sublattice'',
$$
\partial X^{\vee}:=\{\chi^{\vee}\in X^{\vee}\ |\ 
\chi^{\vee}=\sum_{\sigma\in\Gamma_K}\mu_\sigma^\sigma-\mu_\sigma,
\text{ for some }\mu_\sigma\in X^{\vee}\}.
$$
\par
Can the action of $\Gamma$ on $T$ be realized as
an inner automorphism (defined over $\overline{F}$)? By a theorem of
Chevalley, any automorphism
$\sigma\in\Gamma$ may be factored into an inner automorphism and 
an outer automorphism of $G_2$ contained in
$Aut(X,\Delta)$, where $\Delta$ denotes the set of simple positive
roots of $G_2$ with respect to $T/\overline{F}$. Since $G_2$ has no outer
automorphisms, it follows that the answer to our question is yes.
\par
By the previous paragraph, the action of $\Gamma$ on $T$ must factor
through the absolute Weyl group $W$. The only Cartan which had
a $\kappa$-stable germ which did not vanish identically when restricted
to $T$ was the following case:
\newline
Suppose that the non-trivial action of $\Gamma$ on $T$ is realized
by (a) simple reflections $<s_\alpha,s_\zeta>$, (b) the cyclic group
of order $6$, $<s_\beta s_\alpha>$, or (c) the cyclic group
of order $3$, $<(s_\beta s_\alpha)^2>$. 
In the first case, $T$ splits over a
biquadratic extension $E/F$ and $Gal(E/F)$ acts on the character lattice 
by $<s_\alpha,s_\zeta>$. Moreover, $T=T_G$ is isomorphic over $F$
to $T_H=(T_1\times T_2)/\{\pm 1\}
\hookrightarrow (SL(2,F)\times SL(2,F))/\{\pm 1\}$.
In this case, 
$$
H^1(\Gamma,T)=(H^1(\Gamma,T_1)\times H^1(\Gamma,T_2))/\{\pm 1\}
\cong {\Bbb Z}/2{\Bbb Z},
$$
so the character $\kappa$ must be quadratic. 
\par
In the second case, note $s_\beta s_\alpha$ acts by a counterclockwise
rotation on the lattice $X_\Gamma^{\vee}$:
$$
s_\beta s_\alpha (m\alpha^{\vee}+n\beta^{\vee})=
(2m-n)\alpha^{\vee}+(m-n)\beta^{\vee}.
$$
Since
$$
\sum_{i=0}^5(s_\beta s_\alpha)^i (m\alpha^{\vee}+n\beta^{\vee})=0,
$$
we have $X_\Gamma^{\vee}=X^{\vee}$ (this also remains true if we replace
$s_\beta s_\alpha$ by $(s_\beta s_\alpha)^2$). Since
$$
s_\beta s_\alpha (-2\alpha^{\vee}-\beta^{\vee})-(-2\alpha^{\vee}-\beta^{\vee})
=\alpha^{\vee},
$$
$$
s_\beta s_\alpha (-3\alpha^{\vee}-2\beta^{\vee})-
(-3\alpha^{\vee}-2\beta^{\vee})=
\beta^{\vee},
$$
we have $\partial X^{\vee}=X^{\vee}$. Therefore, $H^1(\Gamma,T)$ is trivial by
Tate-Nakayama duality. 
\par
In the cubic case,
$$
s_\beta s_\alpha (m\alpha^{\vee}+n\beta^{\vee})=
(-3n)\alpha^{\vee}+(m-3n)\beta^{\vee}.
$$
Therefore, in this case,
$$
H^1(\Gamma,T)\cong {\Bbb Z}/3{\Bbb Z},
$$
and $\kappa$ must be a character of order $3$.
\par
All other non-split Cartans of $G_2$ do not contribute to the support
of the $2$-regular and $3$-regular $\kappa$-stable germs of $G_2$.
\vskip .1in
\subheading{A.7. The stability of the $2$-regular and $3$-regular
unipotent orbits for $G_2$}
\vskip .1in
I thank M. Assem for explaining to me the technique of this section.
\par
Using the appendices above, we find that
$$
{\frak m}:=cent_{{\frak g}_2}(h_{\alpha +\beta})
={\frak t}+F\cdot X_{3\alpha +\beta}+F\cdot X_{-3\alpha -\beta},
$$
where ${\frak t}$ denotes the Cartan of ${\frak g}_2$. Let
$M:=Cent_{G_2}(h_{\alpha +\beta})$ and let
$$
{\frak g}_2(i):=\{X\in{\frak g}_2\ |\ [h_{\alpha +\beta},X]=iX\}.
$$
We compute
$$
{\frak g}_2(2)=F\cdot X_{\alpha +\beta}.
$$
The map $m\longmapsto Ad(m)X_{\alpha +\beta}$ defines a submersion
$M\rightarrow {\frak g}_2(2)$ \cite{SpSt, ch III, \S 4.14}. By a lemma of
Rao \cite{R}, the $G(F)$-orbit of $X_{\alpha +\beta}$ is
$Ad(K(V(X_{\alpha +\beta})+{\frak n}(2)))$, where
$$
V(X_{\alpha +\beta})=\{ad(X)X_{\alpha +\beta}\ |\ X\in {\frak m}\},
$$
$K$ is a maximal compact subgroup of $G_2(F)$, and 
$$
{\frak n}(2):=\sum_{i >2}{\frak g}_2(i).
$$
The fact that $dim\ {\frak g}_2(2)=1$ gives us that orbit 
of $M$ in ${\frak g}_2(2)$ is onto,
so by Rao's result the $2$-regular unipotent orbit is stable.
\par
A similar calculation shows that the $3$-regular orbit is stable.
(In fact, replacing $h_{\alpha +\beta}$ and $X_{\alpha +\beta}$
by $h_{\beta}$ and $X_{\beta}$, we compute
$$
{\frak g}_2(2)=F\cdot X_{\beta},
$$
and the identical argument holds.)
\par
A Lie triple for the subregular orbit is
$$
X_+=X_\beta+X_{2\alpha +\beta},
X_-=X_{-\beta}+X_{-2\alpha -\beta},
H=h_{\beta}+h_{2\alpha +\beta},
$$
and 
$$
{\frak g}_2(2)=
F\cdot X_{\alpha +\beta}+F\cdot X_{2\alpha +\beta}
+F\cdot X_{3\alpha +\beta}+F\cdot X_{\beta}.
$$
Direct matrix calculations show that the subregular orbit is
not stable. M. Assem has shown that there are 7 $F$-classes in
the stable class, $p\not= 2,3$.

\enddocument
\subheading{A.6. Weighted Dynkin diagrams for $G_2$}
\vskip .1in
Let $e_+,\ e_-, \ f$ be elements of ${\frak g}_2$ satisfying
$[e_+,e_-]=f,[f,e_-]=-2e_-$, and $[f,e_+]=2e_+$. Of course,
this triple gives rise to a 3-dimensional subalgebra ${\frak g}$
of ${\frak g}_2$. Call two such subalgebras ${\frak g}$,
${\frak g}'$ {\bf equivalent} if there is an automorphism
$\phi\in Aut({\frak g}_2)$ for which ${\frak g}'=\phi({\frak g})$.
Following \cite{D, \S 8}, we may attach to
such a subalgebra ${\frak g}$ a weighted diagram depending only on the 
equivalence class of ${\frak g}$. They are constructed as follows:
let
$$
f:= \pmatrix 5/3&\ &\ &\ &\ &\ &\ &\ \\
                         \ &4/3&\ &\ &\ &\ &\ &\ \\
                         \ &\ &1/3&\ &\ &\ &\ &\ \\
                         \ &\ &\ &0&\ &\ &\ &\ \\
                         \ &\ &\ &\ &0&\ &\ &\ \\
                         \ &\ &\ &\ &\ &-1/3&\ &\ \\
                         \ &\ &\ &\ &\ &\ &-4/3&\ \\
                         \ &\ &\ &\ &\ &\ &\ &-5/3
                         \endpmatrix.
$$                                   
Then $f={5\over 3}h_\alpha +3h_\beta$. In general, if
$e_+=c_\alpha X_\alpha +c_\beta X_\beta$ then 
$e_-=c_\alpha X_{-\alpha} +c_\beta X_{-\beta}$ and
$f=2c_\alpha^2 h_\alpha +2c_\beta^2 h_\beta$. From this one can compute 
$e_+$ and $e_-$. The weighted diagram has two nodes connected by
a triple bond, the $\alpha$ node having weight $(f,h_\alpha)=2$ and the 
$\beta$ node having weight $(f,h_\beta)=2$. 
It corresponds to the regular unipotent class.
\par
Let
$$
f:= \pmatrix 2/3&\ &\ &\ &\ &\ &\ &\ \\
                         \ &1/3&\ &\ &\ &\ &\ &\ \\
                         \ &\ &1/3&\ &\ &\ &\ &\ \\
                         \ &\ &\ &0&\ &\ &\ &\ \\
                         \ &\ &\ &\ &0&\ &\ &\ \\
                         \ &\ &\ &\ &\ &-1/3&\ &\ \\
                         \ &\ &\ &\ &\ &\ &-1/3&\ \\
                         \ &\ &\ &\ &\ &\ &\ &-2/3
                         \endpmatrix.
$$                                   
Then $f={2\over 3}h_\alpha +h_\beta$. From this one can compute 
$e_+$ and $e_-$. The weighted diagram has two nodes connected by
a triple bond, the $\alpha$ node having weight $2$ and the 
$\beta$ node having weight $0$. It corresponds to the subregular unipotent
class.
\par
Let
$$
f:= \pmatrix 1/2&\ &\ &\ &\ &\ &\ &\ \\
                         \ &1/2&\ &\ &\ &\ &\ &\ \\
                         \ &\ &0&\ &\ &\ &\ &\ \\
                         \ &\ &\ &0&\ &\ &\ &\ \\
                         \ &\ &\ &\ &0&\ &\ &\ \\
                         \ &\ &\ &\ &\ &0&\ &\ \\
                         \ &\ &\ &\ &\ &\ &-1/2&\ \\
                         \ &\ &\ &\ &\ &\ &\ &-1/2
                         \endpmatrix.
$$                                   
Then $f={1\over 2}h_\alpha +h_\beta$. From this one can compute 
$e_+$ and $e_-$. The weighted diagram has two nodes connected by
a triple bond, the $\alpha$ node having weight $0$ and the 
$\beta$ node having weight $1$. It corresponds to the 2-regular unipotent
class.
\par
Let
$$
f:= \pmatrix 1/3&\ &\ &\ &\ &\ &\ &\ \\
                         \ &1/6&\ &\ &\ &\ &\ &\ \\
                         \ &\ &1/6&\ &\ &\ &\ &\ \\
                         \ &\ &\ &0&\ &\ &\ &\ \\
                         \ &\ &\ &\ &0&\ &\ &\ \\
                         \ &\ &\ &\ &\ &-1/6&\ &\ \\
                         \ &\ &\ &\ &\ &\ &-1/6&\ \\
                         \ &\ &\ &\ &\ &\ &\ &-1/3
                         \endpmatrix.
$$                                   
Then $f={1\over 3}h_\alpha +{1\over 2}h_\beta$. From this one can compute 
$e_+$ and $e_-$. The weighted diagram has two nodes connected by
a triple bond, the $\alpha$ node having weight $1$ and the 
$\beta$ node having weight $0$. It corresponds to the 3-regular unipotent
class.
\vskip .1in



\enddocument
